\paragraph{Why a trigger has a reach.} The crop can only dry the soil at the sensor as fast as the soil delivers water to the roots. That delivery scales with the conductivity, which falls by many orders of magnitude as the soil dries. In a soil with a narrow pore-size distribution (large $n$, sands) the conductivity collapses within a narrow band of suction just past the air entry, so the crop stops drying the soil there; whatever trigger is set beyond that band is never read. In a soil with a broad distribution (small $n$, silty clay loams) the conductivity fades slowly over many decades of suction and the crop can still dry the soil at very high suctions.
\paragraph{Why a deeper sensor does not help.} Below the sensor depth of the experiments the root density is smaller, so the soil dries even less, and in sands the deeper soil is also wetted by drainage from above. A deeper sensor therefore reads lower suctions, not higher ones.
\paragraph{Why scaling by the reach collapses the curves only in part.} The reach sets where irrigation ends, because it is fixed by the conductivity. How irrigation fades before the reach depends on how often the soil reaches a given suction in a season, which depends on the water held between field capacity and that suction, a retention property that differs from the conductivity property that sets the reach.
\paragraph{Why less than half of an irrigation millimetre is transpired.} Transpiration is limited by the energy available to the canopy and by the canopy cover. Water applied to a wet surface or to a canopy that does not cover the ground evaporates from the soil, from the ponded surface and from intercepted water, using the same energy; water applied beyond the storage of the root zone drains.
\paragraph{Why overhead irrigation adds no transpiration.} Water intercepted by the canopy evaporates with the energy that would otherwise drive transpiration, so the canopy transpires less while it is wet. The intercepted evaporation replaces transpiration rather than adding to it, and the marginal transpiration is nil or slightly negative.
\paragraph{Why the dose scale does not depend on the soil.} The water that a dose can lose to drainage is the part that the crop and the surface do not use before it redistributes below the roots. The crop uses about one day of evapotranspiration in that time whatever the soil, so the dose scale is set by the atmosphere and the crop and not by the soil.
\paragraph{Why the plateau depends on the soil.} Once a dose is large, the share that drains is the share of the excess that the soil passes on before the crop uses it, a race between drainage and uptake. Coarse sand passes water fast (high conductivity at field capacity), fine soils hold it, so the plateau is highest in coarse sand.
\paragraph{Why early-season irrigation is more productive.} Early in the season the canopy is growing and its demand rises, so applied water is transpired; late in the season the cereal ripens and is harvested, and the same water is evaporated from bare soil or drains.

